\section{Đặt vấn đề và Động lực}

\subsection{Hiện tượng Nhiễu loạn (Interference) và SUTVA}

Suy luận nhân quả truyền thống dựa trên giả định \textbf{không có nhiễu loạn giữa các đơn vị (Stable Unit Treatment Value Assumption --- SUTVA)}: kết quả của một đơn vị không bị ảnh hưởng bởi đối xử của đơn vị khác. Tuy nhiên, giả định này sụp đổ trong nhiều hệ thống thực tế:

\begin{itemize}
    \item \textbf{Hệ thống gợi ý:} Tương tác của một người dùng làm thay đổi nội dung gợi ý cho những người dùng sau (trạng thái chung: điểm phổ biến, embedding, lịch sử xem).
    \item \textbf{Thị trường / Sàn thương mại điện tử:} Khuyến mãi cho một nhóm người dùng làm thay đổi nhu cầu, giá cả, thời gian chờ và hàng tồn kho (trạng thái chung: lượng tồn kho, giá hiện tại, tín hiệu cầu).
    \item \textbf{Gọi xe công nghệ:} Đối xử với một số hành khách làm thay đổi lượng tài xế trống cho những người khác (trạng thái chung: số tài xế trống, giá linh động, thời gian chờ).
\end{itemize}

Ngay cả \textit{thí nghiệm ngẫu nhiên (randomized experiments)} cũng có thể bị chệch khi một đơn vị được đối xử làm thay đổi điều kiện thị trường chung, ảnh hưởng đến các đơn vị sau.

\subsection{Khái niệm Nhiễu loạn Trạng thái chung (Shared-State Interference)}

Bài báo đề xuất một dạng nhiễu loạn có cấu trúc, được gọi là \textbf{Shared-State Interference (SSI)}: nhiễu loạn không tùy tiện mà chảy qua một trạng thái chung quan sát được \(H_t\), thứ tóm gọn toàn bộ tác động lan tỏa (spillover) trong quá khứ.

\(H_t\) có thể là lượng hàng tồn kho, giá cả, thời gian chờ, điểm phổ biến, hay đầu ra của thuật toán gợi ý --- bất kỳ cơ chế nào trung chuyển tác động lan tỏa giữa các đơn vị theo thời gian. Các đơn vị vừa \textit{bị ảnh hưởng} bởi trạng thái chung, vừa \textit{góp phần làm thay đổi} nó.

Cấu trúc phụ thuộc của SSI được tóm tắt như sau:
\[
    W_{t-1} \;\longrightarrow\; H_t \;\longrightarrow\; Y_t
\]
Quá khứ chỉ ảnh hưởng tương lai thông qua trạng thái chung \(H_t\), không có mũi tên ngang nối trực tiếp giữa các đơn vị.

\subsection{Hạn chế của nghiên cứu trước và Đóng góp của bài báo}

Các nghiên cứu trước đây về nhiễu loạn thường dựa vào mô hình \textbf{tham số cố định (parametric)}, hoặc giả định đặc thù cho từng bài toán cụ thể, dẫn đến dễ bị chệch khi mô hình hóa sai và khó mở rộng sang bối cảnh khác.

Bài báo của Hays và Raghavan (2025) đóng góp ba điểm chính:

\begin{enumerate}
    \item Đưa ra \textbf{khung lý thuyết tổng quát (General Framework)} cho SSI, không phụ thuộc vào giả định tham số cụ thể.
    \item Mở rộng định lý \textbf{Double Machine Learning (DML) cho SSI} --- sử dụng các mô hình học máy linh hoạt, semiparametric, vẫn đảm bảo suy luận thống kê hợp lệ trên dữ liệu phụ thuộc.
    \item Cung cấp \textbf{bộ ước lượng phương sai nhất quán (Consistent Variance Estimator)} để xây dựng khoảng tin cậy chính xác ngay cả khi các quan sát có tương quan theo thời gian.
\end{enumerate}

Điểm mấu chốt: kết hợp được tính linh hoạt của học máy với suy luận thống kê hợp lệ, ngay cả khi dữ liệu phụ thuộc theo chuỗi thời gian.
